\setcounterpageref{ex}{0}
\TaoPhieuKiemTra{18}

\setcounter{ex}{0}
\setcounter{bt}{0}

\textbf{\large I. PHẦN TRẮC NGHIỆM (3,0 điểm)}

{\itshape Khoanh tròn vào chữ cái đặt trước câu trả lời đúng. Mỗi câu đúng được 0,5 điểm.}

\begin{ex} % Câu 1
Số tròn chục lớn nhất có hai chữ số là
\choice
{$99$}
{$100$}
{\True $90$}
{$80$}
\loigiai{
Số tròn chục lớn nhất có hai chữ số là $90$.\\
\textbf{Đáp án đúng: C}
}
\end{ex}

\begin{ex} % Câu 2
Trong phép trừ, nếu số bị trừ bằng số trừ thì hiệu bằng
\choice
{\True $0$}
{$1$}
{Chính số đó}
{Không xác định}
\loigiai{
Khi hai số bằng nhau trừ cho nhau thì hiệu bằng $0$.\\
\textbf{Đáp án đúng: A}
}
\end{ex}

\begin{ex} % Câu 3
Con gà nặng $3\text{ kg}$, con ngỗng nặng $7\text{ kg}$. Cả hai con cân nặng tất cả là
\choice
{$9\text{ kg}$}
{$11\text{ kg}$}
{$4\text{ kg}$}
{\True $10\text{ kg}$}
\loigiai{
$3 + 7 = 10\text{ kg}$.\\
\textbf{Đáp án đúng: D}
}
\end{ex}

\begin{ex} % Câu 4
Sợi dây thừng được căng thẳng tắp cho ta hình ảnh của
\choice
{Đường cong}
{\True Đoạn thẳng}
{Đường gấp khúc}
{Hình tròn}
\loigiai{
Sợi dây kéo căng tạo thành hình ảnh đoạn thẳng.\\
\textbf{Đáp án đúng: B}
}
\end{ex}

\begin{ex} % Câu 5
Điền dấu thích hợp: $7 + 8 \dots 9 + 6$.
\choice
{$>$}
{$<$}
{Không so sánh được}
{\True $=$}
\loigiai{
$7 + 8 = 15$; $9 + 6 = 15$. Vì $15 = 15$ nên điền dấu $=$.\\
\textbf{Đáp án đúng: D}
}
\end{ex}

\begin{ex} % Câu 6
Tính: $15 - 8 + 4 = \dots$
\choice
{$10$}
{$12$}
{\True $11$}
{$9$}
\loigiai{
$15 - 8 = 7$; $7 + 4 = 11$.\\
\textbf{Đáp án đúng: C}
}
\end{ex}

\vspace{4mm}
\textbf{\large II. PHẦN TỰ LUẬN (7,0 điểm)}

{\itshape Thí sinh trình bày chi tiết bài làm vào giấy kiểm tra.}

\begin{bt}[1,0 điểm]
Tính nhẩm:
\begin{multicols}{2}
\begin{enumerate}[a)]
  \item $8 + 7$
  \item $7 + 8$
  \item $15 - 7$
  \item $17 - 8$
\end{enumerate}
\end{multicols}
\loigiai{
\begin{itemize}
  \item $8 + 7 = 15$
  \item $7 + 8 = 15$
  \item $15 - 7 = 8$
  \item $17 - 8 = 9$
\end{itemize}
}
\end{bt}

\begin{bt}[1,5 điểm]
Đặt tính rồi tính:
\begin{multicols}{3}
\begin{enumerate}[a)]
  \item $29 + 36$
  \item $91 - 35$
  \item $52 + 40$
\end{enumerate}
\end{multicols}
\loigiai{
\begin{itemize}
  \item $29 + 36 = 65$
  \item $91 - 35 = 56$
  \item $52 + 40 = 92$
\end{itemize}
}
\end{bt}

\begin{bt}[1,0 điểm]
Tính giá trị biểu thức kèm đơn vị đo:
\begin{enumerate}[a)]
  \item $34\text{ kg} + 9\text{ kg} - 9\text{ kg}$
  \item $43\text{ }l - 14\text{ }l + 8\text{ }l$
\end{enumerate}
\loigiai{
\begin{enumerate}[a)]
  \item $34\text{ kg} + 9\text{ kg} - 9\text{ kg} = 43\text{ kg} - 9\text{ kg} = 34\text{ kg}$.
  \item $43\text{ }l - 14\text{ }l + 8\text{ }l = 29\text{ }l + 8\text{ }l = 37\text{ }l$.
\end{enumerate}
}
\end{bt}

\begin{bt}[1,5 điểm]
Bạn Quân có $33$ ngôi sao giấy. Bạn Sơn có ít hơn bạn Quân $4$ ngôi sao giấy. Hỏi bạn Sơn có bao nhiêu ngôi sao giấy?
\loigiai{
Số ngôi sao giấy bạn Sơn có là:
\[ 33 - 4 = 29 \text{ (ngôi sao giấy)} \]
\textbf{Đáp số:} $29$ ngôi sao giấy.
}
\end{bt}

\begin{bt}[1,0 điểm]
Một thùng chứa đầy $53\text{ }\text{kg}$ gạo nếp. Buổi sáng cửa hàng đã bán ra $14\text{ }\text{kg}$. Hỏi trong thùng còn lại bao nhiêu ki-lô-gam gạo nếp?
\loigiai{
Trong thùng còn lại số ki-lô-gam gạo nếp là:
\[ 53 - 14 = 39\text{ }\text{kg} \]
\textbf{Đáp số:} $39\text{ }\text{kg}$ gạo nếp.
}
\end{bt}

\begin{bt}[1,0 điểm]
Điền tiếp hai số thích hợp vào chỗ chấm để hoàn thành quy luật dãy số: $3, 6, 9, 12, \dots, \dots$
\loigiai{
Quy luật: Mỗi số sau bằng số trước cộng thêm $3$.\\
Số tiếp theo: $12 + 3 = 15$.\\
Số tiếp theo nữa: $15 + 3 = 18$.\\
\textbf{Đáp số:} $15$ và $18$.
}
\end{bt}
